% GENERATED FILE. Edit canonical lessons, then run npm run generate:books.
% canonical-source-hash: fnv1a64:6b885feda38cf1df
% canonical-lessons: JA-C327-reigi, JA-C327-gyougi, JA-C327-shitsurei, JA-C327-meiwaku, JA-C327-ojigi

\chapter{Good Manners}
\label{ch:ja-good-manners}
\hlchaptermodality{327}

\begin{chapteropening}
\textbf{By the end of this chapter:} \emph{I can say what good manners are and why they matter.}

Complete the last lesson of chapter 327: I can say what good manners are and why they matter.
\end{chapteropening}

\section[\texorpdfstring{\ja{れいぎ} (reigi)}{れいぎ (reigi)}]{\ja{れいぎ} (reigi) — courtesy; manners}

\label{lesson:JA-C327-reigi}



\begin{quote}
Before the new one: say the Japanese for a violation, then the Japanese for a punishment.
\end{quote}

\subsection*{You'll want to know: \ja{れいぎ}}
\textbf{\ja{れいぎ}} — \emph{reigi} — \textquotedblleft{}courtesy; manners\textquotedblright{}.

\begin{quote}
\textbf{\ja{こどもに} \ja{れいぎを} \ja{おしえた} \ja{ほうが} \ja{いいです。}} — \emph{kodomo ni reigi o oshieta hō ga ii desu} — \textquotedblleft{}You should teach children good manners.\textquotedblright{}
\end{quote}

\begin{cousinweb}
\textbf{\ja{れいぎ}} (\emph{reigi}) is Sino-Japanese (kango): two Chinese characters read with their on'yomi, readings borrowed from Chinese. They are \emph{rei}, \textquotedblleft{}bow\textquotedblright{} (Middle Chinese \emph{lejX}) and \emph{gi}, \textquotedblleft{}ceremony\textquotedblright{} (\emph{ngje}) (KANJIDIC2; Middle Chinese in Baxter's transcription).
\end{cousinweb}

\begin{grammarlens}[title={What you've built}]
One more word for giving reasons and opinions.
\end{grammarlens}

\subsection*{Your turn}
\begin{itemize}
\raggedright
  \item \emph{Say it:} \emph{reigi}
  \item \emph{Say it:} \emph{reigi}, once more
  \item \emph{Say it:} \emph{kodomo ni reigi o oshieta hō ga ii desu} — You should teach children good manners.
  \item \emph{Recall:} say the Japanese for a violation, then the Japanese for a punishment, then say \emph{reigi} again
\end{itemize}

\begin{tcolorbox}[breakable,colback=teal!4,colframe=teal!35!black,title={Before you move on}]
What does \ja{れいぎ} mean? (\textquotedblleft{}Courtesy; manners\textquotedblright{}.) Say it once more.
\end{tcolorbox}
\section[\texorpdfstring{\ja{ぎょうぎ} (gyōgi)}{ぎょうぎ (gyōgi)}]{\ja{ぎょうぎ} (gyōgi) — behaviour; manners at table}

\label{lesson:JA-C327-gyougi}



\begin{quote}
Before the new one: say the Japanese for a punishment, then the Japanese for courtesy.
\end{quote}

\subsection*{You'll want to know: \ja{ぎょうぎ}}
\textbf{\ja{ぎょうぎ}} — \emph{gyōgi} — \textquotedblleft{}behaviour; manners at table\textquotedblright{}.

\begin{quote}
\textbf{\ja{ぎょうぎが} \ja{いいので、みんなが} \ja{ほめました。}} — \emph{gyōgi ga ii node, minna ga homemashita} — \textquotedblleft{}She behaves well, so everyone praised her.\textquotedblright{}
\end{quote}

\begin{cousinweb}
\textbf{\ja{ぎょうぎ}} (\emph{gyōgi}) is Sino-Japanese (kango): two Chinese characters read with their on'yomi, readings borrowed from Chinese. They are \emph{gyō}, \textquotedblleft{}conduct\textquotedblright{} (Middle Chinese \emph{haengH}) and \emph{gi}, \textquotedblleft{}ceremony\textquotedblright{} (\emph{ngje}) (KANJIDIC2; Middle Chinese in Baxter's transcription).
\end{cousinweb}

\begin{grammarlens}[title={What you've built}]
One more word for giving reasons and opinions.
\end{grammarlens}

\subsection*{Your turn}
\begin{itemize}
\raggedright
  \item \emph{Say it:} \emph{gyōgi}
  \item \emph{Say it:} \emph{gyōgi}, once more
  \item \emph{Say it:} \emph{gyōgi ga ii node, minna ga homemashita} — She behaves well, so everyone praised her.
  \item \emph{Recall:} say the Japanese for a punishment, then the Japanese for courtesy, then say \emph{gyōgi} again
\end{itemize}

\begin{tcolorbox}[breakable,colback=teal!4,colframe=teal!35!black,title={Before you move on}]
What does \ja{ぎょうぎ} mean? (\textquotedblleft{}Behaviour; manners at table\textquotedblright{}.) Say it once more.
\end{tcolorbox}
\section[\texorpdfstring{\ja{しつれい} (shitsurei)}{しつれい (shitsurei)}]{\ja{しつれい} (shitsurei) — rude; impolite}

\label{lesson:JA-C327-shitsurei}



\begin{quote}
Before the new one: say the Japanese for courtesy, then the Japanese for behaviour.
\end{quote}

\subsection*{You'll want to know: \ja{しつれい}}
\textbf{\ja{しつれい}} — \emph{shitsurei} — \textquotedblleft{}rude; impolite\textquotedblright{}.

\begin{quote}
\textbf{\ja{ともだちに} \ja{うそを} \ja{言}\ja{ったら、しつれいです。}} — \emph{tomodachi ni uso o ittara, shitsurei desu} — \textquotedblleft{}Lying to a friend is rude.\textquotedblright{}
\end{quote}

A na-adjective: \ja{な} before a noun.

\begin{cousinweb}
\textbf{\ja{しつれい}} (\emph{shitsurei}) is Sino-Japanese (kango): two Chinese characters read with their on'yomi, readings borrowed from Chinese. They are \emph{shitsu}, \textquotedblleft{}lose\textquotedblright{} and \emph{rei}, \textquotedblleft{}bow\textquotedblright{} (Middle Chinese \emph{lejX}) (KANJIDIC2; Middle Chinese in Baxter's transcription).
\end{cousinweb}

\begin{grammarlens}[title={What you've built}]
One more word for giving reasons and opinions.
\end{grammarlens}

\subsection*{Your turn}
\begin{itemize}
\raggedright
  \item \emph{Say it:} \emph{shitsurei}
  \item \emph{Say it:} \emph{shitsurei}, once more
  \item \emph{Say it:} \emph{tomodachi ni uso o ittara, shitsurei desu} — Lying to a friend is rude.
  \item \emph{Recall:} say the Japanese for courtesy, then the Japanese for behaviour, then say \emph{shitsurei} again
\end{itemize}

\begin{tcolorbox}[breakable,colback=teal!4,colframe=teal!35!black,title={Before you move on}]
What does \ja{しつれい} mean? (\textquotedblleft{}Rude; impolite\textquotedblright{}.) Say it once more.
\end{tcolorbox}
\section[\texorpdfstring{\ja{めいわく} (meiwaku)}{めいわく (meiwaku)}]{\ja{めいわく} (meiwaku) — a nuisance; trouble for others}

\label{lesson:JA-C327-meiwaku}



\begin{quote}
Before the new one: say the Japanese for behaviour, then the Japanese for rude.
\end{quote}

\subsection*{You'll want to know: \ja{めいわく}}
\textbf{\ja{めいわく}} — \emph{meiwaku} — \textquotedblleft{}a nuisance; trouble for others\textquotedblright{}.

\begin{quote}
\textbf{\ja{ばんに} \ja{うたったら、めいわくです。}} — \emph{ban ni utattara, meiwaku desu} — \textquotedblleft{}Singing in the evening is a nuisance.\textquotedblright{}
\end{quote}

A na-adjective: \ja{な} before a noun.

\begin{cousinweb}
\textbf{\ja{めいわく}} (\emph{meiwaku}) is Sino-Japanese (kango): two Chinese characters read with their on'yomi, readings borrowed from Chinese. They are \emph{mei}, \textquotedblleft{}astray\textquotedblright{} (Middle Chinese \emph{mej}) and \emph{waku}, \textquotedblleft{}delusion\textquotedblright{} (\emph{hwok}) (KANJIDIC2; Middle Chinese in Baxter's transcription).
\end{cousinweb}

\begin{grammarlens}[title={What you've built}]
One more word for giving reasons and opinions.
\end{grammarlens}

\subsection*{Your turn}
\begin{itemize}
\raggedright
  \item \emph{Say it:} \emph{meiwaku}
  \item \emph{Say it:} \emph{meiwaku}, once more
  \item \emph{Say it:} \emph{ban ni utattara, meiwaku desu} — Singing in the evening is a nuisance.
  \item \emph{Recall:} say the Japanese for behaviour, then the Japanese for rude, then say \emph{meiwaku} again
\end{itemize}

\begin{tcolorbox}[breakable,colback=teal!4,colframe=teal!35!black,title={Before you move on}]
What does \ja{めいわく} mean? (\textquotedblleft{}A nuisance; trouble for others\textquotedblright{}.) Say it once more.
\end{tcolorbox}
\section[\texorpdfstring{\ja{おじぎ} (ojigi)}{おじぎ (ojigi)}]{\ja{おじぎ} (ojigi) — a bow (as a greeting)}

\label{lesson:JA-C327-ojigi}



\begin{quote}
Before the new one: say the Japanese for rude, then the Japanese for a nuisance.
\end{quote}

\subsection*{You'll want to know: \ja{おじぎ}}
\textbf{\ja{おじぎ}} — \emph{ojigi} — \textquotedblleft{}a bow (as a greeting)\textquotedblright{}.

\begin{quote}
\textbf{\ja{日本人}\ja{は} \ja{おれいの} \ja{とき、おじぎを} \ja{します。}} — \emph{nihonjin wa orei no toki, ojigi o shimasu} — \textquotedblleft{}Japanese people bow when they thank someone.\textquotedblright{}
\end{quote}

\begin{cousinweb}
\textbf{\ja{おじぎ}} (\emph{ojigi}) is built from the native polite prefix \emph{o-} and the Sino-Japanese \emph{jigi}: two Chinese characters read with their on'yomi, readings borrowed from Chinese, \emph{ji}, \textquotedblleft{}word\textquotedblright{} (Middle Chinese \emph{zi}) and \emph{gi}, \textquotedblleft{}ceremony\textquotedblright{} (\emph{ngje}) (KANJIDIC2; JMdict; Middle Chinese in Baxter's transcription).
\end{cousinweb}

\begin{grammarlens}[title={What you've built}]
One more word for giving reasons and opinions.
\end{grammarlens}

\subsection*{Your turn}
\begin{itemize}
\raggedright
  \item \emph{Say it:} \emph{ojigi}
  \item \emph{Say it:} \emph{ojigi}, once more
  \item \emph{Say it:} \emph{nihonjin wa orei no toki, ojigi o shimasu} — Japanese people bow when they thank someone.
  \item \emph{Recall:} say the Japanese for rude, then the Japanese for a nuisance, then say \emph{ojigi} again
\end{itemize}

\begin{tcolorbox}[breakable,colback=teal!4,colframe=teal!35!black,title={Before you move on}]
What does \ja{おじぎ} mean? (\textquotedblleft{}A bow (as a greeting)\textquotedblright{}.) Say it once more.
\end{tcolorbox}
